\chapter{Introduzione}

\section{Contesto e motivazione}
La quantit\`a di dati raccolti da organizzazioni pubbliche e private cresce continuamente. Ospedali, centri di ricerca, aziende e servizi digitali producono informazioni che possono essere utili per costruire modelli statistici, calcolare indicatori o conoscere l'andamento di un fenomeno. Molti di questi dati sono per\`o sensibili. Un dato medico, una misura economica o una preferenza personale non dovrebbe essere consegnata in chiaro a ogni soggetto che partecipa all'elaborazione.

La soluzione pi\`u immediata consiste nel cifrare i dati quando vengono salvati o trasmessi. La cifratura tradizionale protegge bene queste due fasi, ma normalmente richiede di decifrare le informazioni prima di elaborarle. Il server che esegue il calcolo torna quindi a vedere i valori originali. Questo passaggio crea un punto delicato: per ottenere un risultato bisogna concedere fiducia all'infrastruttura che esegue il calcolo.

La cifratura omomorfica affronta il problema da una prospettiva diversa. Essa permette di eseguire alcune operazioni direttamente sui dati cifrati. Il server riceve ciphertext, cio\`e rappresentazioni non leggibili, applica somme o moltiplicazioni e produce un nuovo ciphertext. Dopo la decifratura, il risultato corrisponde al calcolo che sarebbe stato eseguito sui dati in chiaro. Il server pu\`o quindi contribuire all'elaborazione senza conoscere necessariamente gli input.

Questa propriet\`a \`e interessante, ma non rende il calcolo cifrato gratuito o automaticamente adatto a ogni applicazione. Le operazioni omomorfiche sono pi\`u lente delle operazioni ordinarie, richiedono pi\`u memoria e dipendono da parametri che devono essere scelti con attenzione. Schemi diversi rappresentano tipi di dato diversi e offrono compromessi differenti tra precisione, capacit\`a di calcolo e costo.

\section{Obiettivo della tesi}
L'obiettivo di questa tesi \`e studiare in modo progressivo come le tecniche omomorfiche possano essere usate per elaborare e aggregare dati proteggendo maggiormente i contributi individuali. Il lavoro non si limita a descrivere gli schemi da un punto di vista teorico. La parte centrale consiste nel collegare le loro caratteristiche a misure sperimentali e, successivamente, a due applicazioni concrete.

Le domande che guidano il lavoro sono semplici:
\begin{itemize}
  \item quali tipi di dato possono essere rappresentati dai principali schemi omomorfici;
  \item quali parametri incidono maggiormente sul costo delle operazioni;
  \item quale differenza pratica esiste tra sommare e moltiplicare dati cifrati;
  \item come usare queste tecniche in un modello di regressione lineare;
  \item come costruire un protocollo di aggregazione ispirato alle VDAF, nel quale il risultato collettivo possa essere recuperato senza trattare in chiaro ogni contributo.
\end{itemize}

La tesi considera tre schemi disponibili nella libreria OpenFHE: BFV, BGV e CKKS. BFV e BGV sono adatti a operazioni esatte su interi rappresentati in uno spazio modulare. CKKS \`e invece progettato per calcoli approssimati su valori reali o complessi ed \`e quindi utile per molte elaborazioni statistiche e di machine learning \cite{openfhe,ckks}.

\section{Percorso seguito}
Il lavoro \`e iniziato dallo studio tecnologico degli schemi omomorfici. Prima di scegliere uno schema per un'applicazione, \`e infatti necessario capire che cosa possa rappresentare, quali operazioni supporti e come cambino i ciphertext durante il calcolo. Particolare attenzione \`e stata dedicata alla dimensione dell'anello, al numero di slot disponibili, al modulo del plaintext e alla profondit\`a moltiplicativa.

Il secondo passo \`e stato la costruzione di una suite di benchmark basata su OpenFHE. Gli esperimenti misurano il tempo di creazione del contesto, la generazione delle chiavi, la cifratura, la decifratura, la somma e la moltiplicazione omomorfiche. Per CKKS sono stati misurati anche la generazione delle chiavi di bootstrapping e il bootstrapping stesso. I test variano la dimensione dell'anello e la profondit\`a richiesta, in modo da osservare come aumentino tempo e memoria.

I benchmark non servono a dichiarare uno schema vincitore in ogni situazione. Il loro scopo \`e mostrare che la scelta dipende dal problema. Se i dati sono conteggi interi e si desiderano risultati esatti, BFV o BGV sono candidati naturali. Se si elaborano medie, coefficienti o misure non intere, CKKS offre una rappresentazione pi\`u adatta, accettando un piccolo errore numerico.

Dopo il confronto sperimentale, il lavoro applica CKKS a un caso di machine learning semplice e interpretabile: la regressione lineare. Il caso di studio usa dati sanitari e stima la durata di un ricovero. La regressione \`e stata scelta perch\'e permette di osservare con chiarezza il passaggio dai dati cifrati al modello e perch\'e l'inferenza richiede principalmente somme e moltiplicazioni.

L'ultima parte si sposta dal machine learning all'aggregazione verificabile. Viene presentato un protocollo ispirato a VDAF Prio3, ma reinterpretato tramite BGV. I client cifrano report interi, l'aggregatore verifica sotto cifratura che la loro rappresentazione sia valida e somma i contributi accettati. Il collector decifra il totale finale. L'obiettivo non \`e riprodurre integralmente Prio3, ma studiare una costruzione pi\`u semplice che mantenga l'idea di apprendere un aggregato senza esporre direttamente ogni report.

\section{Contributi principali}
I contributi del lavoro possono essere riassunti in quattro punti.

\begin{enumerate}
  \item Una descrizione accessibile delle caratteristiche operative di BFV, BGV e CKKS, con una distinzione tra limiti dello schema, limiti della libreria e limiti della configurazione scelta.
  \item Una suite sperimentale comune per confrontare le principali operazioni omomorfiche al variare della dimensione dell'anello e della profondit\`a moltiplicativa.
  \item Un'applicazione di CKKS alla regressione lineare in uno scenario sanitario, con due possibili distribuzioni dei ruoli di calcolo e decifratura.
  \item Una costruzione di aggregazione VDAF-like basata su BGV, con codifica binaria, validazione omomorfica e decifratura del solo risultato aggregato.
\end{enumerate}

Il filo comune \`e il passaggio dalla singola operazione crittografica a un sistema completo. La scelta dello schema influenza il formato dei dati; i parametri influenzano costo e capacit\`a; l'architettura stabilisce quali attori possano cifrare, calcolare e decifrare. La privacy non dipende quindi da un solo algoritmo, ma dall'insieme di queste decisioni.

\section{Conclusioni}
La cifratura omomorfica permette di ripensare il rapporto tra dati e calcolo: il soggetto che elabora non deve necessariamente vedere ci\`o che elabora. Questa possibilit\`a apre scenari utili, ma introduce costi e vincoli che devono essere misurati e compresi.

La tesi segue quindi un percorso intenzionalmente graduale. Parte dagli strumenti di base, li misura, li applica a un modello statistico e infine li inserisce in un protocollo di aggregazione. In questo modo teoria, benchmark e applicazioni non restano parti separate, ma descrivono fasi successive dello stesso problema.
